\documentclass[11pt,a4paper]{article}

% ============================================================
% PACOTES
% ============================================================

\usepackage[utf8]{inputenc}
\usepackage[T1]{fontenc}
\usepackage[portuguese]{babel}
\usepackage{lmodern}

\usepackage[
a4paper,
top=1.8cm,
bottom=1.8cm,
left=1.7cm,
right=1.7cm
]{geometry}

\usepackage{amsmath,amssymb}
\usepackage{array}
\usepackage{booktabs}
\usepackage{tabularx}
\usepackage{multirow}
\usepackage{multicol}
\usepackage{enumitem}
\usepackage{float}
\usepackage{xcolor}
\usepackage{graphicx}

% ============================================================
% SIUNITX -- conforme solicitado
% ============================================================

\usepackage{siunitx}
\sisetup{output-decimal-marker={,}}

% Fórmulas químicas
\usepackage[version=4]{mhchem}

% TikZ
\usepackage{tikz}
\usetikzlibrary{
	arrows.meta,
	positioning,
	calc,
	shapes.geometric,
	decorations.pathmorphing,
	backgrounds,
	fit
}

% Gráficos
\usepackage{pgfplots}
\pgfplotsset{compat=1.18}

% Caixas
\usepackage[most]{tcolorbox}

% Cabeçalhos
\usepackage{fancyhdr}
\usepackage{lastpage}

\usepackage{hyperref}

% ============================================================
% CORES
% ============================================================

\definecolor{eublue}{HTML}{145A86}
\definecolor{eudark}{HTML}{17324D}
\definecolor{eulightblue}{HTML}{E4F2FA}

\definecolor{eugreen}{HTML}{3F7D58}
\definecolor{eulightgreen}{HTML}{E8F4EC}

\definecolor{euorange}{HTML}{E28A2B}
\definecolor{eulightorange}{HTML}{FFF0DE}

\definecolor{eured}{HTML}{B33A3A}
\definecolor{eupurple}{HTML}{67518A}

\definecolor{eugray}{HTML}{F1F3F5}
\definecolor{euyellow}{HTML}{F3CF55}

\definecolor{metalcolor}{HTML}{DDECF6}
\definecolor{nonmetalcolor}{HTML}{E4F3E6}
\definecolor{noblegascolor}{HTML}{F4E2F5}
\definecolor{metalloidcolor}{HTML}{FFF0C9}

% ============================================================
% CONFIGURAÇÕES
% ============================================================

\setlength{\parindent}{0pt}
\setlength{\parskip}{4pt}

\renewcommand{\arraystretch}{1.35}

\hypersetup{
	colorlinks=true,
	linkcolor=eublue,
	urlcolor=eublue
}

% ============================================================
% CABEÇALHO E RODAPÉ
% ============================================================

\pagestyle{fancy}
\fancyhf{}

\fancyhead[L]{
	\textbf{\textcolor{eublue}{EUFISICA exercícios}}
}

\fancyhead[R]{
	\textcolor{eudark}{Física e Química A -- 10.º ano}
}

\fancyfoot[L]{Prof. José Gonçalves}
\fancyfoot[C]{EUFISICA}
\fancyfoot[R]{Página \thepage\ de \pageref{LastPage}}

\renewcommand{\headrulewidth}{0.6pt}
\renewcommand{\footrulewidth}{0.4pt}

% ============================================================
% CAIXAS
% ============================================================

\newcounter{exercicio}

\newenvironment{exercicio}[1]
{
	\refstepcounter{exercicio}
	\begin{tcolorbox}[
		enhanced,
		breakable,
		colback=white,
		colframe=eublue,
		boxrule=0.8pt,
		arc=2mm,
		left=3mm,
		right=3mm,
		top=2mm,
		bottom=2mm,
		title={
			\textbf{Exercício \theexercicio}
			\hspace{1mm}--\hspace{1mm}#1
		},
		coltitle=white,
		colbacktitle=eublue,
		fonttitle=\bfseries
		]
	}
	{
	\end{tcolorbox}
}

\newtcolorbox{recorda}{
	enhanced,
	breakable,
	colback=eulightblue,
	colframe=eublue,
	title=\textbf{Recorda},
	coltitle=eudark,
	colbacktitle=eulightblue,
	fonttitle=\bfseries,
	arc=2mm
}

\newtcolorbox{dados}{
	enhanced,
	breakable,
	colback=eugray,
	colframe=eudark,
	title=\textbf{Dados},
	fonttitle=\bfseries,
	arc=2mm
}

\newtcolorbox{stem}{
	enhanced,
	breakable,
	colback=eulightgreen,
	colframe=eugreen,
	title=\textbf{STEM | Ciência, Tecnologia e Mundo Real},
	coltitle=eugreen,
	colbacktitle=eulightgreen,
	fonttitle=\bfseries,
	arc=2mm
}

\newtcolorbox{desafio}{
	enhanced,
	breakable,
	colback=eulightorange,
	colframe=euorange,
	title=\textbf{Desafio},
	coltitle=euorange,
	colbacktitle=eulightorange,
	fonttitle=\bfseries,
	arc=2mm
}

% ============================================================
% LINHAS DE RESPOSTA
% ============================================================

\newcommand{\linha}{
	\vspace{3mm}
	\noindent\makebox[\linewidth]{\dotfill}
}

\newcommand{\linhas}[1]{
	\par
	\foreach \i in {1,...,#1}{
		\vspace{3.5mm}
		\noindent\makebox[\linewidth]{\dotfill}\par
	}
}

% ============================================================
% ELEMENTO DA TABELA PERIÓDICA
% ============================================================

\tikzset{
	elem/.style={
		draw=eudark,
		minimum width=0.68cm,
		minimum height=0.72cm,
		inner sep=1pt,
		align=center,
		font=\scriptsize
	}
}

\newcommand{\Elemento}[5]{%
	\node[
	elem,
	fill=#5
	] at ({(#1-1)*0.72},{-(#2-1)*0.78})
	{
		{\tiny #3}\\[-1mm]
		{\bfseries #4}
	};
}

% ============================================================
% DOCUMENTO
% ============================================================

\begin{document}
	
	% ============================================================
	% TÍTULO
	% ============================================================
	
	\begin{tcolorbox}[
		enhanced,
		colback=eudark,
		colframe=eudark,
		arc=3mm,
		boxrule=0pt,
		left=5mm,
		right=5mm,
		top=5mm,
		bottom=5mm
		]
		
		\begin{center}
			
			{\Huge\bfseries\color{white}
				EUFISICA exercícios}
			
			\vspace{2mm}
			
			{\Large\bfseries\color{white}
				Tabela Periódica e Estrutura Atómica}
			
			\vspace{2mm}
			
			{\large\color{eulightblue}
				Organização • Propriedades periódicas • Distribuição eletrónica}
			
			{\large\color{eulightblue}
				Espetros • Energia • Quantidade de matéria}
			
			\vspace{3mm}
			
			{\large\color{white}
				\textbf{Física e Química A -- 10.º ano}}
			
			\vspace{2mm}
			
			{\color{white}
				\textbf{Prof. José Gonçalves}}
			
		\end{center}
		
	\end{tcolorbox}
	
	\vspace{3mm}
	
	\begin{tabularx}{\textwidth}{|X|X|X|}
		\hline
		\textbf{Nome:}
		&
		\textbf{Turma:}
		&
		\textbf{Data:}
		\\[6mm]
		\hline
	\end{tabularx}
	
	\vspace{5mm}
	
	% ============================================================
	% OBJETIVOS
	% ============================================================
	
	\begin{tcolorbox}[
		colback=eugray,
		colframe=eudark,
		title=\textbf{Objetivos de aprendizagem}
		]
		
		No final desta ficha deverás ser capaz de:
		
		\begin{itemize}[leftmargin=6mm,itemsep=1mm]
			
			\item compreender a organização da Tabela Periódica;
			
			\item relacionar número atómico, distribuição eletrónica,
			período e grupo;
			
			\item identificar eletrões de valência;
			
			\item prever tendências do raio atómico e da energia de ionização;
			
			\item relacionar propriedades periódicas com a estrutura eletrónica;
			
			\item interpretar energias de remoção eletrónica e espetros
			fotoeletrónicos;
			
			\item interpretar espetros de emissão e de absorção;
			
			\item calcular a energia de fotões e de níveis eletrónicos;
			
			\item relacionar massa, massa molar, quantidade de matéria
			e número de entidades;
			
			\item integrar diferentes conhecimentos de estrutura atómica
			na identificação de elementos químicos.
			
		\end{itemize}
		
	\end{tcolorbox}
	
	% ============================================================
	% DADOS
	% ============================================================
	
	\begin{dados}
		
		Sempre que necessário, utiliza:
		
		\[
		N_A=\SI{6.022e23}{\per\mole}
		\]
		
		\[
		h=\SI{6.626e-34}{\joule\second}
		\]
		
		\[
		c=\SI{3.00e8}{\meter\per\second}
		\]
		
		\[
		1\,\mathrm{eV}
		=
		\SI{1.602e-19}{\joule}
		\]
		
		\[
		n=\frac{m}{M}
		\]
		
		\[
		N=nN_A
		\]
		
		\[
		E_{\text{fotão}}=h\nu
		\]
		
		\[
		c=\lambda\nu
		\]
		
		\[
		E_{\text{fotão}}
		=
		\frac{hc}{\lambda}
		\]
		
		Para o átomo de hidrogénio:
		
		\[
		\boxed{
			E_n=
			-\frac{2,18\times10^{-18}}{n^2}
			\ \mathrm{J}
		}
		\]
		
		\[
		\Delta E=E_f-E_i
		\]
		
	\end{dados}
	
	% ============================================================
	% 1. ORGANIZAÇÃO TP
	% ============================================================
	
	\section*{\textcolor{eublue}{1. Organização da Tabela Periódica}}
	
	\begin{recorda}
		
		Os elementos químicos encontram-se organizados na Tabela Periódica
		por ordem crescente de \textbf{número atómico}, \(Z\).
		
		\[
		\boxed{
			Z=\text{número de protões}
		}
		\]
		
		Num átomo neutro:
		
		\[
		\boxed{
			n(e^-)=Z
		}
		\]
		
		As linhas horizontais denominam-se \textbf{períodos}.
		
		As colunas verticais denominam-se \textbf{grupos}.
		
		Para os elementos representativos:
		
		\begin{itemize}
			
			\item o período relaciona-se com o maior valor do número quântico
			principal \(n\) ocupado;
			
			\item elementos do mesmo grupo apresentam distribuições eletrónicas
			de valência semelhantes;
			
			\item essa semelhança explica muitas das suas propriedades químicas.
			
		\end{itemize}
		
	\end{recorda}
	
	% ============================================================
	% TABELA PERIÓDICA SIMPLIFICADA
	% ============================================================
	
	\begin{figure}[H]
		\centering
		
		\resizebox{\textwidth}{!}{
			\begin{tikzpicture}
				
				% números dos grupos
				\foreach \g in {1,...,18}{
					\node[
					font=\scriptsize\bfseries,
					text=eudark
					] at ({(\g-1)*0.72},0.72)
					{\g};
				}
				
				\node[
				font=\scriptsize\bfseries,
				rotate=90
				] at (-1.0,-1.2)
				{Período};
				
				% períodos
				\foreach \p in {1,...,4}{
					\node[
					font=\scriptsize\bfseries
					] at (-0.65,{-(\p-1)*0.78})
					{\p};
				}
				
				% período 1
				\Elemento{1}{1}{1}{H}{nonmetalcolor}
				\Elemento{18}{1}{2}{He}{noblegascolor}
				
				% período 2
				\Elemento{1}{2}{3}{Li}{metalcolor}
				\Elemento{2}{2}{4}{Be}{metalcolor}
				\Elemento{13}{2}{5}{B}{metalloidcolor}
				\Elemento{14}{2}{6}{C}{nonmetalcolor}
				\Elemento{15}{2}{7}{N}{nonmetalcolor}
				\Elemento{16}{2}{8}{O}{nonmetalcolor}
				\Elemento{17}{2}{9}{F}{nonmetalcolor}
				\Elemento{18}{2}{10}{Ne}{noblegascolor}
				
				% período 3
				\Elemento{1}{3}{11}{Na}{metalcolor}
				\Elemento{2}{3}{12}{Mg}{metalcolor}
				\Elemento{13}{3}{13}{Al}{metalcolor}
				\Elemento{14}{3}{14}{Si}{metalloidcolor}
				\Elemento{15}{3}{15}{P}{nonmetalcolor}
				\Elemento{16}{3}{16}{S}{nonmetalcolor}
				\Elemento{17}{3}{17}{Cl}{nonmetalcolor}
				\Elemento{18}{3}{18}{Ar}{noblegascolor}
				
				% período 4
				\Elemento{1}{4}{19}{K}{metalcolor}
				\Elemento{2}{4}{20}{Ca}{metalcolor}
				
				% legenda
				\node[
				draw=eudark,
				fill=metalcolor,
				minimum width=1.3cm,
				minimum height=0.45cm
				] at (2.0,-3.5) {};
				
				\node[
				anchor=west,
				font=\scriptsize
				] at (2.75,-3.5) {metal};
				
				\node[
				draw=eudark,
				fill=metalloidcolor,
				minimum width=1.3cm,
				minimum height=0.45cm
				] at (5.5,-3.5) {};
				
				\node[
				anchor=west,
				font=\scriptsize
				] at (6.25,-3.5) {semimetal};
				
				\node[
				draw=eudark,
				fill=nonmetalcolor,
				minimum width=1.3cm,
				minimum height=0.45cm
				] at (9.0,-3.5) {};
				
				\node[
				anchor=west,
				font=\scriptsize
				] at (9.75,-3.5) {não metal};
				
				\node[
				draw=eudark,
				fill=noblegascolor,
				minimum width=1.3cm,
				minimum height=0.45cm
				] at (12.7,-3.5) {};
				
				\node[
				anchor=west,
				font=\scriptsize
				] at (13.45,-3.5) {gás nobre};
				
			\end{tikzpicture}
		}
		
		\caption{Representação simplificada dos primeiros 20 elementos da Tabela
			Periódica.}
	\end{figure}
	
	% ============================================================
	
	\begin{exercicio}{Conhecer a Tabela Periódica}
		
		Utiliza a Tabela Periódica representada anteriormente.
		
		Completa:
		
		\begin{center}
			
			\small
			
			\begin{tabular}{|c|c|c|c|c|}
				\hline
				
				\textbf{Elemento}
				&
				\(\boldsymbol{Z}\)
				&
				\textbf{Período}
				&
				\textbf{Grupo}
				&
				\textbf{Metal/não metal}
				\\
				\hline
				
				\ce{Na}
				&
				&
				&
				&
				\\[4mm]
				\hline
				
				\ce{Mg}
				&
				&
				&
				&
				\\[4mm]
				\hline
				
				\ce{O}
				&
				&
				&
				&
				\\[4mm]
				\hline
				
				\ce{Cl}
				&
				&
				&
				&
				\\[4mm]
				\hline
				
				\ce{Ne}
				&
				&
				&
				&
				\\[4mm]
				\hline
				
				\ce{Ca}
				&
				&
				&
				&
				\\[4mm]
				\hline
				
			\end{tabular}
			
		\end{center}
		
	\end{exercicio}
	
	% ============================================================
	
	\begin{exercicio}{A posição revela informação}
		
		Responde sem recorrer diretamente ao símbolo do elemento.
		
		\begin{enumerate}[label=\alph*)]
			
			\item Um elemento encontra-se no grupo 1 e período 3.
			
			Identifica-o.
			
			\linha
			
			\item Um elemento encontra-se no grupo 17 e período 2.
			
			Identifica-o.
			
			\linha
			
			\item Um elemento encontra-se no grupo 18 e período 3.
			
			Identifica-o.
			
			\linha
			
			\item Um elemento possui número atómico \(Z=12\).
			
			Indica:
			
			\begin{itemize}
				\item símbolo;
				\item período;
				\item grupo.
			\end{itemize}
			
			\linhas{2}
			
			\item Qual dos seguintes elementos pertence ao mesmo grupo do sódio?
			
			\[
			\ce{Li},
			\qquad
			\ce{Mg},
			\qquad
			\ce{K},
			\qquad
			\ce{Cl}
			\]
			
			\linha
			
		\end{enumerate}
		
	\end{exercicio}
	
	% ============================================================
	% CONFIGURAÇÃO ELETRÓNICA
	% ============================================================
	
	\section*{\textcolor{eublue}{2. Distribuição eletrónica e posição na Tabela Periódica}}
	
	\begin{recorda}
		
		Para os primeiros elementos, a ordem de preenchimento dos subníveis é:
		
		\[
		1s
		\rightarrow
		2s
		\rightarrow
		2p
		\rightarrow
		3s
		\rightarrow
		3p
		\rightarrow
		4s
		\]
		
		Por exemplo:
		
		\[
		\ce{Na}:
		\qquad
		1s^2\,2s^2\,2p^6\,3s^1
		\]
		
		O maior valor de \(n\) é 3.
		
		Logo, o sódio encontra-se no:
		
		\[
		\boxed{\text{3.º período}}
		\]
		
		A configuração de valência é:
		
		\[
		3s^1
		\]
		
		o que está de acordo com a sua posição no grupo 1.
		
	\end{recorda}
	
	% ============================================================
	
	\begin{exercicio}{Da configuração eletrónica à Tabela Periódica}
		
		Completa a tabela.
		
		\begin{center}
			
			\small
			
			\begin{tabular}{|l|c|c|c|c|}
				\hline
				
				\textbf{Distribuição eletrónica}
				&
				\(\boldsymbol{Z}\)
				&
				\textbf{Elemento}
				&
				\textbf{Período}
				&
				\textbf{Grupo}
				\\
				\hline
				
				\(1s^2\,2s^1\)
				&
				&
				&
				&
				\\[4mm]
				\hline
				
				\(1s^2\,2s^2\,2p^4\)
				&
				&
				&
				&
				\\[4mm]
				\hline
				
				\(1s^2\,2s^2\,2p^6\,3s^1\)
				&
				&
				&
				&
				\\[4mm]
				\hline
				
				\(1s^2\,2s^2\,2p^6\,3s^2\)
				&
				&
				&
				&
				\\[4mm]
				\hline
				
				\(1s^2\,2s^2\,2p^6\,3s^2\,3p^5\)
				&
				&
				&
				&
				\\[4mm]
				\hline
				
				\(1s^2\,2s^2\,2p^6\,3s^2\,3p^6\)
				&
				&
				&
				&
				\\[4mm]
				\hline
				
			\end{tabular}
			
		\end{center}
		
	\end{exercicio}
	
	% ============================================================
	
	\begin{exercicio}{Eletrões de valência}
		
		Completa.
		
		\begin{center}
			
			\begin{tabular}{|c|c|c|c|}
				\hline
				
				\textbf{Elemento}
				&
				\textbf{Configuração de valência}
				&
				\textbf{Eletrões de valência}
				&
				\textbf{Grupo}
				\\
				\hline
				
				\ce{Li}
				&
				&
				&
				\\[4mm]
				\hline
				
				\ce{Mg}
				&
				&
				&
				\\[4mm]
				\hline
				
				\ce{Al}
				&
				&
				&
				\\[4mm]
				\hline
				
				\ce{O}
				&
				&
				&
				\\[4mm]
				\hline
				
				\ce{Cl}
				&
				&
				&
				\\[4mm]
				\hline
				
				\ce{Ne}
				&
				&
				&
				\\[4mm]
				\hline
				
			\end{tabular}
			
		\end{center}
		
		Que relação identificas entre o número de eletrões de valência e o grupo
		nos elementos representativos?
		
		\linhas{3}
		
	\end{exercicio}
	
	% ============================================================
	% ORBITAIS
	% ============================================================
	
	\begin{exercicio}{Do diagrama de orbitais ao elemento}
		
		Considera o seguinte diagrama de orbitais.
		
		\begin{center}
			
			\begin{tikzpicture}[scale=0.9]
				
				% 1s
				\node at (-0.5,0.45) {\(1s\)};
				\draw (0,0) rectangle (0.9,0.9);
				\draw[-{Latex},thick,eublue]
				(0.25,0.15)--(0.25,0.75);
				\draw[{Latex}-,thick,eured]
				(0.65,0.15)--(0.65,0.75);
				
				% 2s
				\node at (1.45,0.45) {\(2s\)};
				\draw (1.9,0) rectangle (2.8,0.9);
				\draw[-{Latex},thick,eublue]
				(2.15,0.15)--(2.15,0.75);
				\draw[{Latex}-,thick,eured]
				(2.55,0.15)--(2.55,0.75);
				
				% 2p
				\node at (3.4,0.45) {\(2p\)};
				
				\foreach \x in {3.9,4.8,5.7}{
					\draw (\x,0) rectangle ({\x+0.9},0.9);
					
					\draw[-{Latex},thick,eublue]
					({\x+0.25},0.15)--({\x+0.25},0.75);
					
					\draw[{Latex}-,thick,eured]
					({\x+0.65},0.15)--({\x+0.65},0.75);
				}
				
				% 3s
				\node at (7.1,0.45) {\(3s\)};
				\draw (7.6,0) rectangle (8.5,0.9);
				
				\draw[-{Latex},thick,eublue]
				(7.85,0.15)--(7.85,0.75);
				
				\draw[{Latex}-,thick,eured]
				(8.25,0.15)--(8.25,0.75);
				
			\end{tikzpicture}
			
		\end{center}
		
		\begin{enumerate}[label=\alph*)]
			
			\item Quantos eletrões apresenta o átomo?
			
			\linha
			
			\item Qual é o seu número atómico?
			
			\linha
			
			\item Escreve a distribuição eletrónica.
			
			\linhas{1}
			
			\item Identifica o elemento.
			
			\linha
			
			\item Indica o período.
			
			\linha
			
			\item Indica o grupo.
			
			\linha
			
			\item Quantos eletrões de valência possui?
			
			\linha
			
			\item Prevê o ião monoatómico mais provável que este elemento poderá
			formar.
			
			\linhas{2}
			
		\end{enumerate}
		
	\end{exercicio}
	
	% ============================================================
	% IÕES E PERIODICIDADE
	% ============================================================
	
	\begin{exercicio}{Átomos e iões}
		
		Completa a tabela.
		
		\begin{center}
			
			\small
			
			\begin{tabular}{|c|c|c|c|c|}
				\hline
				
				\textbf{Espécie}
				&
				\textbf{N.º eletrões}
				&
				\textbf{Distribuição eletrónica}
				&
				\textbf{Período do átomo}
				&
				\textbf{Grupo}
				\\
				\hline
				
				\(\ce{Na}\)
				&
				&
				&
				&
				\\[4mm]
				\hline
				
				\(\ce{Na+}\)
				&
				&
				&
				&
				\\[4mm]
				\hline
				
				\(\ce{Mg^{2+}}\)
				&
				&
				&
				&
				\\[4mm]
				\hline
				
				\(\ce{O^{2-}}\)
				&
				&
				&
				&
				\\[4mm]
				\hline
				
				\(\ce{Cl-}\)
				&
				&
				&
				&
				\\[4mm]
				\hline
				
			\end{tabular}
			
		\end{center}
		
		\begin{enumerate}[label=\alph*)]
			
			\item Quais das espécies são isoeletrónicas?
			
			\linhas{2}
			
			\item Explica por que motivo a posição de um elemento na Tabela Periódica
			é determinada pelo átomo neutro e não pelo ião formado.
			
			\linhas{3}
			
		\end{enumerate}
		
	\end{exercicio}
	
	% ============================================================
	% 3. PROPRIEDADES PERIÓDICAS
	% ============================================================
	
	\newpage
	
	\section*{\textcolor{eublue}{3. Propriedades periódicas}}
	
	\begin{recorda}
		
		Ao longo de um período, o número de protões aumenta.
		
		Os eletrões adicionados encontram-se aproximadamente na mesma camada
		eletrónica principal.
		
		De forma geral:
		
		\[
		\boxed{
			\text{ao longo de um período, o raio atómico diminui}
		}
		\]
		
		e:
		
		\[
		\boxed{
			\text{a energia de ionização tende a aumentar}
		}
		\]
		
		Ao descer num grupo, aumenta o número de camadas eletrónicas ocupadas.
		
		Assim, de forma geral:
		
		\[
		\boxed{
			\text{o raio atómico aumenta ao descer num grupo}
		}
		\]
		
		e:
		
		\[
		\boxed{
			\text{a primeira energia de ionização diminui}
		}
		\]
		
	\end{recorda}
	
	% ============================================================
	% TENDÊNCIAS
	% ============================================================
	
	\begin{figure}[H]
		\centering
		
		\begin{tikzpicture}[scale=0.88]
			
			% grelha estilizada
			\foreach \x in {0,...,5}{
				\foreach \y in {0,...,3}{
					\draw[eudark!50]
					(\x,\y) rectangle +(0.75,0.6);
				}
			}
			
			\node[
			font=\bfseries,
			text=eublue
			] at (1.8,3.3)
			{Raio atómico aumenta};
			
			\draw[
			-{Latex[length=4mm]},
			ultra thick,
			eublue
			]
			(4.2,2.8)--(0.2,2.8);
			
			\draw[
			-{Latex[length=4mm]},
			ultra thick,
			eublue
			]
			(-0.45,3.0)--(-0.45,0.1);
			
			\node[
			font=\bfseries,
			text=eured
			] at (3.2,-1.0)
			{Energia de ionização aumenta};
			
			\draw[
			-{Latex[length=4mm]},
			ultra thick,
			eured
			]
			(0.3,-0.45)--(4.5,-0.45);
			
			\draw[
			-{Latex[length=4mm]},
			ultra thick,
			eured
			]
			(5.0,0.1)--(5.0,2.9);
			
		\end{tikzpicture}
		
		\caption{Tendências gerais do raio atómico e da primeira energia de
			ionização na Tabela Periódica.}
	\end{figure}
	
	% ============================================================
	
	\begin{exercicio}{Comparar raios atómicos}
		
		Sem consultar valores tabelados, indica em cada par o átomo que,
		em geral, deverá apresentar maior raio atómico.
		
		Justifica com base na posição na Tabela Periódica.
		
		\begin{enumerate}[label=\alph*)]
			
			\item \(\ce{Li}\) ou \(\ce{Na}\)
			
			\linhas{2}
			
			\item \(\ce{Na}\) ou \(\ce{Cl}\)
			
			\linhas{2}
			
			\item \(\ce{Mg}\) ou \(\ce{Ca}\)
			
			\linhas{2}
			
			\item \(\ce{O}\) ou \(\ce{F}\)
			
			\linhas{2}
			
			\item Ordena por ordem crescente de raio atómico:
			
			\[
			\ce{Na},
			\qquad
			\ce{Mg},
			\qquad
			\ce{Al},
			\qquad
			\ce{Cl}
			\]
			
			\linhas{2}
			
		\end{enumerate}
		
	\end{exercicio}
	
	% ============================================================
	% GRÁFICO ENERGIA IONIZAÇÃO PERÍODO 3
	% ============================================================
	
	\begin{figure}[H]
		\centering
		
		\begin{tikzpicture}
			
			\begin{axis}[
				width=13cm,
				height=7cm,
				xlabel={Elemento},
				ylabel={Primeira energia de ionização / \(\mathrm{kJ\,mol^{-1}}\)},
				symbolic x coords={Na,Mg,Al,Si,P,S,Cl,Ar},
				xtick=data,
				ymin=0,
				ymax=1700,
				grid=major,
				grid style={dashed,eugray},
				mark=*,
				mark size=2.5pt,
				thick,
				axis line style={eudark},
				label style={font=\small\bfseries}
				]
				
				\addplot[
				very thick,
				mark=*
				]
				coordinates{
					(Na,496)
					(Mg,738)
					(Al,578)
					(Si,787)
					(P,1012)
					(S,1000)
					(Cl,1251)
					(Ar,1521)
				};
				
			\end{axis}
			
		\end{tikzpicture}
		
		\caption{Primeira energia de ionização dos elementos do 3.º período.}
	\end{figure}
	
	% ============================================================
	
	\begin{exercicio}{Interpretar uma propriedade periódica}
		
		Analisa o gráfico anterior.
		
		\begin{enumerate}[label=\alph*)]
			
			\item Descreve a tendência geral da primeira energia de ionização ao
			longo do 3.º período.
			
			\linhas{2}
			
			\item Qual dos elementos representados apresenta a menor primeira energia
			de ionização?
			
			\linha
			
			\item Qual apresenta a maior?
			
			\linha
			
			\item Explica a tendência geral observada tendo em conta:
			
			\begin{itemize}
				\item o aumento da carga nuclear;
				\item a distância dos eletrões de valência ao núcleo.
			\end{itemize}
			
			\linhas{4}
			
			\item Verifica-se uma diminuição entre Mg e Al.
			
			Observa:
			
			\[
			\ce{Mg}:
			1s^2\,2s^2\,2p^6\,3s^2
			\]
			
			\[
			\ce{Al}:
			1s^2\,2s^2\,2p^6\,3s^2\,3p^1
			\]
			
			Usa as distribuições eletrónicas para propor uma explicação para a
			diferença observada.
			
			\linhas{4}
			
			\item Existe também uma pequena diminuição de P para S.
			
			Compara:
			
			\[
			\ce{P}:3p^3
			\]
			
			e
			
			\[
			\ce{S}:3p^4.
			\]
			
			Relaciona a diferença com o emparelhamento de eletrões nos orbitais
			\(p\).
			
			\linhas{4}
			
		\end{enumerate}
		
	\end{exercicio}
	
	% ============================================================
	% GRUPO 1
	% ============================================================
	
	\begin{figure}[H]
		\centering
		
		\begin{tikzpicture}
			
			\begin{axis}[
				width=11cm,
				height=6.5cm,
				ybar,
				bar width=18pt,
				ylabel={1.ª energia de ionização / \(\mathrm{kJ\,mol^{-1}}\)},
				xlabel={Elemento do grupo 1},
				symbolic x coords={Li,Na,K,Rb,Cs},
				xtick=data,
				ymin=0,
				ymax=600,
				grid=major,
				grid style={dashed,eugray},
				nodes near coords,
				axis line style={eudark}
				]
				
				\addplot[
				fill=eublue!60,
				draw=eublue
				]
				coordinates{
					(Li,520)
					(Na,496)
					(K,419)
					(Rb,403)
					(Cs,376)
				};
				
			\end{axis}
			
		\end{tikzpicture}
		
		\caption{Primeira energia de ionização de alguns elementos do grupo 1.}
	\end{figure}
	
	% ============================================================
	
	\begin{exercicio}{Descer num grupo}
		
		Analisa o gráfico anterior.
		
		\begin{enumerate}[label=\alph*)]
			
			\item Qual é a tendência da primeira energia de ionização ao descer no
			grupo 1?
			
			\linhas{2}
			
			\item Como varia o raio atómico no mesmo sentido?
			
			\linhas{2}
			
			\item Explica por que motivo o eletrão de valência do césio é, em geral,
			mais fácil de remover do que o eletrão de valência do lítio.
			
			\linhas{4}
			
			\item Qual dos elementos representados esperarias que formasse mais
			facilmente um ião \(+1\)?
			
			Justifica.
			
			\linhas{3}
			
		\end{enumerate}
		
	\end{exercicio}
	
	% ============================================================
	% 4. NÚMERO ATÓMICO / MASSA ATÓMICA
	% ============================================================
	
	\section*{\textcolor{eublue}{4. Número atómico, isótopos e massa atómica relativa}}
	
	\begin{recorda}
		
		Isótopos são átomos do mesmo elemento que possuem:
		
		\[
		\boxed{
			\text{o mesmo número de protões}
		}
		\]
		
		mas diferente número de neutrões.
		
		Por isso, apresentam o mesmo \(Z\), mas diferente \(A\).
		
		A massa atómica relativa de um elemento é uma média ponderada das
		massas dos seus isótopos naturais.
		
	\end{recorda}
	
	% ============================================================
	
	\begin{exercicio}{O cloro na Tabela Periódica}
		
		O cloro natural pode ser aproximado por uma mistura de:
		
		\begin{center}
			
			\begin{tabular}{ccc}
				\toprule
				\textbf{Isótopo}
				&
				\textbf{Massa isotópica}
				&
				\textbf{Abundância}
				\\
				\midrule
				
				\({}^{35}\ce{Cl}\)
				&
				34,969
				&
				\(75,77\,\%\)
				\\
				
				\({}^{37}\ce{Cl}\)
				&
				36,966
				&
				\(24,23\,\%\)
				\\
				
				\bottomrule
			\end{tabular}
			
		\end{center}
		
		\begin{enumerate}[label=\alph*)]
			
			\item Quantos protões possui cada isótopo?
			
			\linha
			
			\item Quantos neutrões possui aproximadamente
			\({}^{35}_{17}\ce{Cl}\)?
			
			\linha
			
			\item E \({}^{37}_{17}\ce{Cl}\)?
			
			\linha
			
			\item Calcula a massa atómica relativa do cloro.
			
			\linhas{3}
			
			\item Compara o resultado com o valor indicado habitualmente na Tabela
			Periódica.
			
			\linhas{2}
			
			\item Explica por que motivo os dois isótopos ocupam a mesma posição na
			Tabela Periódica.
			
			\linhas{3}
			
		\end{enumerate}
		
	\end{exercicio}
	
	% ============================================================
	% 5. QUANTIDADE DE MATÉRIA
	% ============================================================
	
	\section*{\textcolor{eublue}{5. Tabela Periódica, massa molar e quantidade de matéria}}
	
	\begin{recorda}
		
		A massa molar de uma substância pode ser determinada utilizando as
		massas atómicas relativas dos elementos que a constituem.
		
		\[
		\boxed{
			n=\frac{m}{M}
		}
		\]
		
		e:
		
		\[
		\boxed{
			N=nN_A
		}
		\]
		
	\end{recorda}
	
	% ============================================================
	
	\begin{exercicio}{Da Tabela Periódica para a massa molar}
		
		Utiliza os seguintes valores:
		
		\[
		A_r(\ce{H})=1,008
		\]
		
		\[
		A_r(\ce{C})=12,01
		\]
		
		\[
		A_r(\ce{O})=16,00
		\]
		
		\[
		A_r(\ce{Na})=22,99
		\]
		
		\[
		A_r(\ce{Mg})=24,31
		\]
		
		\[
		A_r(\ce{Cl})=35,45
		\]
		
		\[
		A_r(\ce{Ca})=40,08
		\]
		
		Calcula as massas molares de:
		
		\begin{multicols}{2}
			
			\begin{enumerate}[label=\alph*)]
				
				\item \(\ce{H2O}\)
				
				\item \(\ce{CO2}\)
				
				\item \(\ce{NaCl}\)
				
				\item \(\ce{MgCl2}\)
				
				\item \(\ce{CaCO3}\)
				
				\item \(\ce{Na2CO3}\)
				
			\end{enumerate}
			
		\end{multicols}
		
		\linhas{5}
		
	\end{exercicio}
	
	% ============================================================
	
	\begin{exercicio}{Uma amostra de magnésio}
		
		Uma fita de magnésio apresenta massa:
		
		\[
		m=\SI{4.86}{g}
		\]
		
		Considera:
		
		\[
		M(\ce{Mg})=\SI{24.31}{g\per\mole}
		\]
		
		\begin{enumerate}[label=\alph*)]
			
			\item Determina a quantidade de matéria de magnésio.
			
			\linhas{2}
			
			\item Determina o número de átomos de magnésio.
			
			\linhas{2}
			
			\item Quantos protões existem, no total, nessa amostra?
			
			Recorda que:
			
			\[
			Z(\ce{Mg})=12.
			\]
			
			\linhas{3}
			
			\item Quantos eletrões existem na amostra, admitindo átomos neutros?
			
			\linhas{2}
			
			\item A que grupo e período pertence o magnésio?
			
			\linhas{2}
			
		\end{enumerate}
		
	\end{exercicio}
	
	% ============================================================
	% 6. ESPETROS
	% ============================================================
	
	\newpage
	
	\section*{\textcolor{eublue}{6. Elementos químicos e espetros atómicos}}
	
	\begin{recorda}
		
		Os eletrões dos átomos só podem ocupar determinados estados de energia.
		
		Quando um eletrão passa de um estado de maior energia para outro de menor
		energia, pode ocorrer emissão de um fotão.
		
		\[
		\boxed{
			E_{\text{fotão}}=|\Delta E|
		}
		\]
		
		Cada elemento possui um conjunto próprio de estados de energia e,
		consequentemente, um espetro característico.
		
	\end{recorda}
	
	% ============================================================
	% ESPECTROS
	% ============================================================
	
	\begin{figure}[H]
		\centering
		
		\begin{tikzpicture}[x=0.018cm,y=1cm]
			
			% H
			\node[
			anchor=east,
			font=\bfseries
			] at (370,3.7)
			{\ce{H}};
			
			\fill[black]
			(380,3.3) rectangle (700,4.1);
			
			\draw[violet,very thick]
			(410,3.3)--(410,4.1);
			
			\draw[blue!70!violet,very thick]
			(434,3.3)--(434,4.1);
			
			\draw[cyan!70!blue,very thick]
			(486,3.3)--(486,4.1);
			
			\draw[red,very thick]
			(656,3.3)--(656,4.1);
			
			% Na
			\node[
			anchor=east,
			font=\bfseries
			] at (370,2.4)
			{\ce{Na}};
			
			\fill[black]
			(380,2.0) rectangle (700,2.8);
			
			\draw[yellow!80!orange,very thick]
			(589,2.0)--(589,2.8);
			
			\draw[yellow!80!orange,very thick]
			(589.6,2.0)--(589.6,2.8);
			
			% He
			\node[
			anchor=east,
			font=\bfseries
			] at (370,1.1)
			{\ce{He}};
			
			\fill[black]
			(380,0.7) rectangle (700,1.5);
			
			\foreach \x in {447,502,588,668}{
				\draw[euorange,very thick]
				(\x,0.7)--(\x,1.5);
			}
			
			% X
			\node[
			anchor=east,
			font=\bfseries
			] at (370,-0.2)
			{X};
			
			\fill[black]
			(380,-0.6) rectangle (700,0.2);
			
			\foreach \x in {410,434,486,589,656}{
				\draw[eugreen!70!yellow,very thick]
				(\x,-0.6)--(\x,0.2);
			}
			
			% eixo
			\draw[-{Latex},thick]
			(380,-1.0)--(710,-1.0);
			
			\foreach \x in {400,450,500,550,600,650,700}{
				\draw
				(\x,-1.08)--(\x,-0.92);
				
				\node[
				below,
				font=\scriptsize
				] at (\x,-1.08)
				{\x};
			}
			
			\node at (540,-1.55)
			{\small Comprimento de onda / nm};
			
		\end{tikzpicture}
		
		\caption{Espetros de emissão esquemáticos de alguns elementos e de uma
			amostra X.}
	\end{figure}
	
	% ============================================================
	
	\begin{exercicio}{A impressão digital dos elementos}
		
		Analisa os espetros anteriores.
		
		\begin{enumerate}[label=\alph*)]
			
			\item Que elementos parecem existir na amostra X?
			
			\linhas{2}
			
			\item Explica por que razão uma risca a \(\SI{589}{nm}\) pode fornecer
			informação sobre a presença de sódio.
			
			\linhas{3}
			
			\item A posição do sódio na Tabela Periódica é:
			
			\[
			\text{grupo }1,\qquad\text{período }3.
			\]
			
			Escreve a sua distribuição eletrónica.
			
			\linhas{1}
			
			\item Quantos eletrões de valência possui?
			
			\linha
			
			\item Qual o ião mais provável formado pelo sódio?
			
			\linha
			
			\item Explica por que motivo o espetro do sódio não pode ser previsto
			apenas sabendo que pertence ao grupo 1.
			
			\linhas{3}
			
		\end{enumerate}
		
	\end{exercicio}
	
	% ============================================================
	
	\begin{exercicio}{Energia de uma risca do sódio}
		
		Uma das riscas intensas do sódio apresenta:
		
		\[
		\lambda=\SI{589}{nm}
		\]
		
		\begin{enumerate}[label=\alph*)]
			
			\item Converte o comprimento de onda para metros.
			
			\linhas{1}
			
			\item Determina a frequência da radiação.
			
			\linhas{2}
			
			\item Calcula a energia de um fotão.
			
			\linhas{2}
			
			\item Calcula a energia correspondente a um mol destes fotões,
			em \(\si{\kilo\joule\per\mole}\).
			
			\linhas{3}
			
			\item A emissão desta radiação corresponde a uma transição eletrónica
			para um nível de maior ou menor energia?
			
			Justifica.
			
			\linhas{2}
			
		\end{enumerate}
		
	\end{exercicio}
	
	% ============================================================
	% 7. HIDROGÉNIO
	% ============================================================
	
	\section*{\textcolor{eublue}{7. O hidrogénio: Tabela Periódica e níveis de energia}}
	
	\begin{recorda}
		
		Para o átomo de hidrogénio:
		
		\[
		\boxed{
			E_n=
			-\frac{2,18\times10^{-18}}{n^2}
			\ \mathrm{J}
		}
		\]
		
		A energia aumenta com \(n\), aproximando-se de zero.
		
		Para:
		
		\[
		n\rightarrow\infty
		\]
		
		tem-se:
		
		\[
		E=0.
		\]
		
		Numa transição:
		
		\[
		\boxed{
			\Delta E=E_f-E_i
		}
		\]
		
		Se:
		
		\[
		\Delta E<0,
		\]
		
		ocorre emissão de energia.
		
		Se:
		
		\[
		\Delta E>0,
		\]
		
		ocorre absorção de energia.
		
	\end{recorda}
	
	% ============================================================
	% NÍVEIS H
	% ============================================================
	
	\begin{figure}[H]
		\centering
		
		\begin{tikzpicture}[x=1cm,y=0.55cm]
			
			\draw[
			-{Latex[length=3mm]},
			thick
			]
			(0,-14.3)--(0,1.2);
			
			\node[
			rotate=90,
			font=\bfseries
			] at (-0.65,-6.5)
			{Energia / eV};
			
			\draw[very thick,eudark]
			(1,-13.6)--(7,-13.6);
			
			\node[left]
			at (1,-13.6)
			{\(n=1\)};
			
			\node[right]
			at (7,-13.6)
			{\(-13,6\)};
			
			\draw[very thick,eudark]
			(1,-3.40)--(7,-3.40);
			
			\node[left]
			at (1,-3.40)
			{\(n=2\)};
			
			\node[right]
			at (7,-3.40)
			{\(-3,40\)};
			
			\draw[very thick,eudark]
			(1,-1.51)--(7,-1.51);
			
			\node[left]
			at (1,-1.51)
			{\(n=3\)};
			
			\node[right]
			at (7,-1.51)
			{\(-1,51\)};
			
			\draw[very thick,eudark]
			(1,-0.85)--(7,-0.85);
			
			\node[left]
			at (1,-0.85)
			{\(n=4\)};
			
			\node[right]
			at (7,-0.85)
			{\(-0,85\)};
			
			\draw[very thick,eudark]
			(1,-0.54)--(7,-0.54);
			
			\node[left]
			at (1,-0.54)
			{\(n=5\)};
			
			\node[right]
			at (7,-0.54)
			{\(-0,54\)};
			
			\draw[
			dashed,
			very thick,
			eured
			]
			(1,0)--(7,0);
			
			\node[left]
			at (1,0)
			{\(n=\infty\)};
			
			\node[right,text=eured]
			at (7,0)
			{\(0\)};
			
			% transições
			\draw[
			-{Latex[length=3mm]},
			very thick,
			eublue
			]
			(2.5,-0.85)--(2.5,-3.40);
			
			\node[
			right,
			text=eublue
			] at (2.55,-2.2)
			{\small A};
			
			\draw[
			-{Latex[length=3mm]},
			very thick,
			eugreen
			]
			(4.5,-13.6)--(4.5,-3.40);
			
			\node[
			right,
			text=eugreen
			] at (4.55,-8.5)
			{\small B};
			
			\draw[
			-{Latex[length=3mm]},
			very thick,
			euorange
			]
			(6,-1.51)--(6,-13.6);
			
			\node[
			right,
			text=euorange
			] at (6.05,-7)
			{\small C};
			
		\end{tikzpicture}
		
		\caption{Estados de energia do átomo de hidrogénio.}
	\end{figure}
	
	% ============================================================
	
	\begin{exercicio}{O elemento mais simples}
		
		O hidrogénio apresenta:
		
		\[
		Z=1.
		\]
		
		\begin{enumerate}[label=\alph*)]
			
			\item Escreve a distribuição eletrónica do átomo de hidrogénio no estado
			fundamental.
			
			\linha
			
			\item Qual o seu período?
			
			\linha
			
			\item Em que grupo se encontra?
			
			\linha
			
			\item O hidrogénio possui uma posição particular na Tabela Periódica.
			Apesar de ser colocado no grupo 1, não é um metal alcalino.
			
			Que informação da sua distribuição eletrónica justifica a colocação
			no grupo 1?
			
			\linhas{2}
			
			\item Determina a energia do nível \(n=1\), em joule.
			
			\linhas{2}
			
			\item Determina a energia do nível \(n=2\).
			
			\linhas{2}
			
		\end{enumerate}
		
	\end{exercicio}
	
	% ============================================================
	
	\begin{exercicio}{Transições no átomo de hidrogénio}
		
		Considera as transições A, B e C representadas no diagrama.
		
		\begin{enumerate}[label=\alph*)]
			
			\item Classifica cada transição como emissão ou absorção.
			
			\linhas{2}
			
			\item Para a transição A:
			
			\[
			n=4\rightarrow n=2
			\]
			
			determina:
			
			\begin{enumerate}[label=\roman*)]
				
				\item \(E_4\);
				
				\item \(E_2\);
				
				\item \(\Delta E\);
				
				\item a energia do fotão;
				
				\item a frequência do fotão;
				
				\item o comprimento de onda.
				
			\end{enumerate}
			
			\linhas{7}
			
			\item Relaciona o comprimento de onda obtido com o espetro de emissão do
			hidrogénio apresentado anteriormente.
			
			\linhas{2}
			
			\item Qual das transições representadas deverá envolver maior quantidade
			de energia?
			
			Justifica sem efetuar cálculos detalhados.
			
			\linhas{3}
			
		\end{enumerate}
		
	\end{exercicio}
	
	% ============================================================
	% 8. FOTOELETRÓNICO
	% ============================================================
	
	\newpage
	
	\section*{\textcolor{eublue}{8. Espetro fotoeletrónico e Tabela Periódica}}
	
	\begin{recorda}
		
		A espetroscopia fotoeletrónica permite estudar as energias de remoção
		dos eletrões.
		
		\[
		\boxed{
			E_{\text{fotão}}
			=
			E_{\text{remoção}}
			+
			E_c
		}
		\]
		
		ou:
		
		\[
		\boxed{
			E_{\text{remoção}}
			=
			E_{\text{fotão}}
			-
			E_c
		}
		\]
		
		A posição dos picos fornece informação sobre a energia de remoção.
		
		A intensidade relativa dos picos está relacionada com o número de
		eletrões existente em cada subnível.
		
	\end{recorda}
	
	% ============================================================
	% PES MG
	% ============================================================
	
	\begin{figure}[H]
		\centering
		
		\begin{tikzpicture}[x=1.45cm,y=0.7cm]
			
			% eixos
			\draw[
			-{Latex[length=3mm]},
			thick
			]
			(0,0)--(8.5,0);
			
			\draw[
			-{Latex[length=3mm]},
			thick
			]
			(0,0)--(0,7.5);
			
			\node[
			rotate=90,
			font=\bfseries
			] at (-0.7,3.7)
			{Intensidade relativa};
			
			\node[
			font=\bfseries
			] at (4.2,-1.2)
			{Energia de remoção};
			
			% sentido maior energia
			\draw[
			-{Latex[length=3mm]},
			thick,
			eured
			]
			(7.4,-0.6)--(0.9,-0.6);
			
			\node[
			text=eured,
			font=\small
			] at (4.2,-0.6)
			{energia de remoção crescente};
			
			% Mg: 1s2, 2s2, 2p6, 3s2
			
			\draw[eublue,very thick]
			(1.1,0)--(1.1,2.0);
			
			\node[above,font=\bfseries]
			at (1.1,2.0)
			{\(1s\)};
			
			\draw[eublue,very thick]
			(3.0,0)--(3.0,2.0);
			
			\node[above,font=\bfseries]
			at (3.0,2.0)
			{\(2s\)};
			
			\draw[eublue,very thick]
			(4.8,0)--(4.8,6.0);
			
			\node[above,font=\bfseries]
			at (4.8,6.0)
			{\(2p\)};
			
			\draw[eublue,very thick]
			(7.0,0)--(7.0,2.0);
			
			\node[above,font=\bfseries]
			at (7.0,2.0)
			{\(3s\)};
			
		\end{tikzpicture}
		
		\caption{Espetro fotoeletrónico esquemático de um elemento X. As
			distâncias horizontais não estão à escala.}
	\end{figure}
	
	% ============================================================
	
	\begin{exercicio}{Identificar um elemento a partir do espetro fotoeletrónico}
		
		O espetro anterior corresponde a um elemento X.
		
		\begin{enumerate}[label=\alph*)]
			
			\item Quantos subníveis eletrónicos estão ocupados?
			
			\linha
			
			\item Com base nas intensidades relativas, determina o número de eletrões
			em:
			
			\[
			1s,\qquad2s,\qquad2p,\qquad3s.
			\]
			
			\linhas{2}
			
			\item Escreve a distribuição eletrónica completa do elemento.
			
			\linhas{1}
			
			\item Determina o número total de eletrões.
			
			\linha
			
			\item Determina \(Z\).
			
			\linha
			
			\item Identifica o elemento na Tabela Periódica.
			
			\linha
			
			\item Indica o período.
			
			\linha
			
			\item Indica o grupo.
			
			\linha
			
			\item Qual dos subníveis apresenta os eletrões mais facilmente removíveis?
			
			\linha
			
			\item Explica por que motivo esses eletrões são mais fáceis de remover.
			
			\linhas{3}
			
		\end{enumerate}
		
	\end{exercicio}
	
	% ============================================================
	% IONIZAÇÕES SUCESSIVAS
	% ============================================================
	
	\begin{exercicio}{Energias de ionização sucessivas}
		
		Para um elemento Y obtiveram-se:
		
		\begin{center}
			
			\begin{tabular}{cc}
				\toprule
				
				\textbf{Eletrão removido}
				&
				\textbf{Energia / \(\si{\kilo\joule\per\mole}\)}
				\\
				\midrule
				
				1.º & 738\\
				2.º & 1451\\
				3.º & 7733\\
				4.º & 10543\\
				5.º & 13630\\
				
				\bottomrule
			\end{tabular}
			
		\end{center}
		
		\begin{enumerate}[label=\alph*)]
			
			\item Entre que remoções se observa o maior aumento de energia?
			
			\linha
			
			\item Quantos eletrões de valência deverá possuir o átomo?
			
			\linha
			
			\item A que grupo da Tabela Periódica deverá pertencer?
			
			\linha
			
			\item Sabendo que pertence ao 3.º período, identifica o elemento.
			
			\linha
			
			\item Escreve a distribuição eletrónica.
			
			\linhas{1}
			
			\item Representa a configuração de valência num diagrama de orbitais.
			
			\vspace{12mm}
			
			\item Explica por que motivo a terceira energia de ionização é muito
			superior à segunda.
			
			\linhas{4}
			
		\end{enumerate}
		
	\end{exercicio}
	
	% ============================================================
	% 9. TABELA PERIÓDICA COMO FERRAMENTA
	% ============================================================
	
	\section*{\textcolor{eublue}{9. A Tabela Periódica como ferramenta de previsão}}
	
	\begin{exercicio}{Quem sou eu?}
		
		Identifica cada elemento a partir das pistas.
		
		\begin{enumerate}[label=\alph*)]
			
			\item Tenho \(Z=8\), pertenço ao 2.º período e possuo seis eletrões de
			valência.
			
			\[
			\boxed{\text{Elemento: }\hspace{3cm}}
			\]
			
			\item A minha distribuição eletrónica termina em \(3s^1\).
			
			\[
			\boxed{\text{Elemento: }\hspace{3cm}}
			\]
			
			\item Sou um gás nobre do 3.º período.
			
			\[
			\boxed{\text{Elemento: }\hspace{3cm}}
			\]
			
			\item Tenho configuração de valência:
			
			\[
			3s^2\,3p^5.
			\]
			
			\[
			\boxed{\text{Elemento: }\hspace{3cm}}
			\]
			
			\item Formo habitualmente um ião \(2+\), pertenço ao 3.º período e a
			minha configuração eletrónica termina em \(3s^2\).
			
			\[
			\boxed{\text{Elemento: }\hspace{3cm}}
			\]
			
			\item Possuo 10 eletrões e sou quimicamente muito pouco reativo.
			
			\[
			\boxed{\text{Elemento: }\hspace{3cm}}
			\]
			
		\end{enumerate}
		
	\end{exercicio}
	
	% ============================================================
	
	\begin{exercicio}{Verdadeiro ou falso}
		
		Classifica cada afirmação como verdadeira (V) ou falsa (F).
		Corrige as falsas.
		
		\begin{enumerate}
			
			\item[\(\square\)]
			Os elementos estão organizados na Tabela Periódica por ordem crescente
			de massa atómica.
			
			\item[\(\square\)]
			Elementos do mesmo grupo possuem sempre o mesmo número total de
			eletrões.
			
			\item[\(\square\)]
			Para elementos representativos, elementos do mesmo grupo apresentam
			configurações de valência semelhantes.
			
			\item[\(\square\)]
			Ao longo de um período, o número atómico aumenta.
			
			\item[\(\square\)]
			Ao descer num grupo, o número de camadas eletrónicas ocupadas aumenta.
			
			\item[\(\square\)]
			A energia de ionização aumenta sempre, sem qualquer exceção, ao longo
			de um período.
			
			\item[\(\square\)]
			Um espetro fotoeletrónico pode fornecer informação sobre a distribuição
			eletrónica.
			
			\item[\(\square\)]
			Todos os elementos de um mesmo grupo apresentam espetros de emissão
			iguais.
			
			\item[\(\square\)]
			A massa molar de um elemento está relacionada com a sua massa atómica
			relativa.
			
		\end{enumerate}
		
		\linhas{5}
		
	\end{exercicio}
	
	% ============================================================
	% PROBLEMA INTEGRADOR
	% ============================================================
	
	\newpage
	
	\section*{\textcolor{euorange}{10. Problema integrador -- Identificar um elemento desconhecido}}
	
	\begin{desafio}
		
		Uma equipa de investigação analisa um elemento metálico desconhecido,
		designado por X.
		
		Os cientistas obtiveram as seguintes informações:
		
		\begin{itemize}
			
			\item o átomo neutro possui 11 eletrões;
			
			\item o espetro fotoeletrónico é compatível com:
			
			\[
			1s^2\,2s^2\,2p^6\,3s^1;
			\]
			
			\item uma das linhas de emissão mais intensas ocorre próximo de:
			
			\[
			\lambda=\SI{589}{nm};
			\]
			
			\item a massa molar é aproximadamente:
			
			\[
			M=\SI{22.99}{g\per\mole};
			\]
			
			\item uma amostra apresenta massa:
			
			\[
			m=\SI{2.30}{g}.
			\]
			
		\end{itemize}
		
		Responde.
		
		\begin{enumerate}[label=\alph*)]
			
			\item Determina o número atómico de X.
			
			\linhas{1}
			
			\item Identifica X.
			
			\linhas{1}
			
			\item Localiza o elemento na Tabela Periódica, indicando grupo e período.
			
			\linhas{2}
			
			\item Quantos eletrões de valência apresenta?
			
			\linha
			
			\item Representa os eletrões de valência num diagrama de orbitais.
			
			\vspace{12mm}
			
			\item Que ião monoatómico é expectável que forme?
			
			Justifica recorrendo à distribuição eletrónica.
			
			\linhas{3}
			
			\item Determina a frequência da radiação de \(\SI{589}{nm}\).
			
			\linhas{2}
			
			\item Determina a energia de um fotão dessa radiação.
			
			\linhas{2}
			
			\item A linha observada resulta de absorção ou de emissão de energia,
			sabendo que se trata de um espetro de emissão?
			
			\linhas{2}
			
			\item Determina a quantidade de matéria existente na amostra de
			\(\SI{2.30}{g}\).
			
			\linhas{2}
			
			\item Determina o número de átomos presentes na amostra.
			
			\linhas{2}
			
			\item Cada átomo possui 11 protões.
			
			Estima o número total de protões existentes na amostra.
			
			\linhas{3}
			
			\item O elemento imediatamente à direita de X no mesmo período apresenta
			maior primeira energia de ionização.
			
			Identifica esse elemento e explica a tendência geral.
			
			\linhas{3}
			
			\item O elemento imediatamente abaixo de X no mesmo grupo apresenta
			maior raio atómico.
			
			Explica porquê.
			
			\linhas{3}
			
		\end{enumerate}
		
	\end{desafio}
	
	% ============================================================
	% STEM
	% ============================================================
	
	\section*{\textcolor{eugreen}{11. Tabela Periódica e mundo real}}
	
	\begin{stem}
		
		\textbf{De que são feitas as estrelas?}
		
		A Tabela Periódica não é apenas uma tabela utilizada em laboratório.
		
		Os elementos que nela aparecem constituem também:
		
		\begin{itemize}
			
			\item estrelas;
			
			\item planetas;
			
			\item atmosferas;
			
			\item seres vivos;
			
			\item materiais tecnológicos;
			
			\item baterias;
			
			\item semicondutores;
			
			\item medicamentos;
			
			\item painéis solares.
			
		\end{itemize}
		
		Através da \textbf{espetroscopia}, é possível identificar elementos
		presentes numa estrela a partir da luz que recebemos.
		
		Através da \textbf{espetroscopia fotoeletrónica}, é possível investigar
		a distribuição dos eletrões num material.
		
		A posição de um elemento na Tabela Periódica permite ainda prever muitas
		das suas propriedades.
		
		Assim, a Tabela Periódica representa uma ligação entre:
		
		\[
		\boxed{
			\text{estrutura atómica}
			\longrightarrow
			\text{propriedades}
			\longrightarrow
			\text{materiais}
			\longrightarrow
			\text{tecnologia}
		}
		\]
		
	\end{stem}
	
	% ============================================================
	
	\begin{exercicio}{Escolher um elemento para uma aplicação}
		
		Uma equipa pretende escolher materiais para diferentes aplicações.
		
		Relaciona cada situação com uma propriedade periódica ou eletrónica
		relevante.
		
		\begin{enumerate}[label=\alph*)]
			
			\item Pretende-se um elemento que perca facilmente um eletrão para
			formar um ião \(+1\).
			
			Que região da Tabela Periódica procurarias?
			
			\linhas{2}
			
			\item Pretende-se um elemento quimicamente pouco reativo.
			
			Que grupo seria um ponto de partida razoável?
			
			\linhas{2}
			
			\item Pretende-se identificar sódio numa amostra distante sem recolher
			fisicamente o material.
			
			Que técnica estudada poderias usar?
			
			\linhas{2}
			
			\item Pretende-se estudar as energias de ligação dos eletrões de um
			sólido.
			
			Que tipo de espetro seria útil?
			
			\linhas{2}
			
			\item Explica por que razão a Tabela Periódica pode ser considerada uma
			ferramenta de previsão e não apenas uma lista de elementos.
			
			\linhas{4}
			
		\end{enumerate}
		
	\end{exercicio}
	
	% ============================================================
	% DESAFIO FINAL
	% ============================================================
	
	\section*{\textcolor{euorange}{12. Desafio final}}
	
	\begin{desafio}
		
		Três elementos A, B e C encontram-se no 3.º período.
		
		Sabe-se que:
		
		\begin{itemize}
			
			\item A possui configuração de valência \(3s^1\);
			
			\item B possui configuração de valência \(3s^2\,3p^4\);
			
			\item C apresenta uma camada de valência completa;
			
			\item a primeira energia de ionização de A é muito inferior à de C.
			
		\end{itemize}
		
		\begin{enumerate}[label=\alph*)]
			
			\item Identifica A.
			
			\linha
			
			\item Identifica B.
			
			\linha
			
			\item Identifica C.
			
			\linha
			
			\item Indica os grupos de A, B e C.
			
			\linhas{2}
			
			\item Ordena A, B e C por ordem crescente de primeira energia de
			ionização.
			
			\linhas{2}
			
			\item Ordena-os, qualitativamente, por ordem crescente de raio atómico.
			
			\linhas{2}
			
			\item Qual deverá formar mais facilmente um catião?
			
			Justifica.
			
			\linhas{3}
			
			\item Qual deverá apresentar menor tendência para reagir?
			
			Justifica.
			
			\linhas{3}
			
			\item Explica como a distribuição eletrónica permite prever estas
			diferenças de comportamento.
			
			\linhas{4}
			
		\end{enumerate}
		
	\end{desafio}
	
	% ============================================================
	% AUTOAVALIAÇÃO
	% ============================================================
	
	\section*{\textcolor{eublue}{13. Autoavaliação}}
	
	\begin{center}
		
		\small
		
		\begin{tabularx}{\textwidth}{
				>{\raggedright\arraybackslash}X
				c c c
			}
			
			\toprule
			
			&
			\textbf{Tenho dúvidas}
			&
			\textbf{Estou a aprender}
			&
			\textbf{Consigo sozinho}
			\\
			
			\midrule
			
			Localizo um elemento pelo seu \(Z\).
			&
			\(\square\)&\(\square\)&\(\square\)
			\\[3mm]
			
			Relaciono configuração eletrónica, período e grupo.
			&
			\(\square\)&\(\square\)&\(\square\)
			\\[3mm]
			
			Identifico eletrões de valência.
			&
			\(\square\)&\(\square\)&\(\square\)
			\\[3mm]
			
			Comparo raios atómicos.
			&
			\(\square\)&\(\square\)&\(\square\)
			\\[3mm]
			
			Interpreto energias de ionização.
			&
			\(\square\)&\(\square\)&\(\square\)
			\\[3mm]
			
			Interpreto um espetro fotoeletrónico.
			&
			\(\square\)&\(\square\)&\(\square\)
			\\[3mm]
			
			Relaciono espetros de emissão com elementos.
			&
			\(\square\)&\(\square\)&\(\square\)
			\\[3mm]
			
			Calculo a energia de um fotão.
			&
			\(\square\)&\(\square\)&\(\square\)
			\\[3mm]
			
			Utilizo a expressão \(E_n\) para o hidrogénio.
			&
			\(\square\)&\(\square\)&\(\square\)
			\\[3mm]
			
			Calculo quantidade de matéria e número de entidades.
			&
			\(\square\)&\(\square\)&\(\square\)
			\\[3mm]
			
			Integro diferentes informações para identificar um elemento.
			&
			\(\square\)&\(\square\)&\(\square\)
			\\
			
			\bottomrule
			
		\end{tabularx}
		
	\end{center}
	
	% ============================================================
	% MAPA CONCEPTUAL FINAL
	% ============================================================
	
	\section*{\textcolor{eublue}{14. Mapa conceptual}}
	
	\begin{center}
		
		\begin{tikzpicture}[
			node distance=1.8cm and 2cm,
			caixa/.style={
				draw=eudark,
				rounded corners,
				minimum width=3.2cm,
				minimum height=1.05cm,
				align=center,
				font=\small\bfseries
			}
			]
			
			\node[
			caixa,
			fill=eulightblue
			] (tp)
			{Tabela\\Periódica};
			
			\node[
			caixa,
			fill=eulightgreen,
			below left=of tp
			] (z)
			{Número\\atómico};
			
			\node[
			caixa,
			fill=eulightorange,
			below=of tp
			] (config)
			{Distribuição\\eletrónica};
			
			\node[
			caixa,
			fill=noblegascolor,
			below right=of tp
			] (prop)
			{Propriedades\\periódicas};
			
			\node[
			caixa,
			fill=eugray,
			below=of z
			] (mol)
			{Quantidade\\de matéria};
			
			\node[
			caixa,
			fill=eugray,
			below=of config
			] (pes)
			{Espetro\\fotoeletrónico};
			
			\node[
			caixa,
			fill=eugray,
			below=of prop
			] (espetros)
			{Espetros de\\emissão/absorção};
			
			\node[
			caixa,
			fill=eulightblue,
			below=of pes
			] (energia)
			{Energia dos\\eletrões};
			
			\node[
			caixa,
			fill=eulightgreen,
			below=of espetros
			] (foton)
			{Energia dos\\fotões};
			
			\draw[-{Latex},thick]
			(tp)--(z);
			
			\draw[-{Latex},thick]
			(tp)--(config);
			
			\draw[-{Latex},thick]
			(tp)--(prop);
			
			\draw[-{Latex},thick]
			(z)--(mol);
			
			\draw[-{Latex},thick]
			(config)--(pes);
			
			\draw[-{Latex},thick]
			(prop)--(espetros);
			
			\draw[-{Latex},thick]
			(pes)--(energia);
			
			\draw[-{Latex},thick]
			(espetros)--(foton);
			
		\end{tikzpicture}
		
	\end{center}
	
	\vfill
	
	% ============================================================
	% RODAPÉ VISUAL
	% ============================================================
	
	\begin{center}
		
		\begin{tikzpicture}
			
			\draw[
			eublue,
			very thick
			]
			(0,0)--(13,0);
			
			\node[
			text=eublue,
			font=\Large\bfseries
			]
			at (6.5,0.55)
			{EUFISICA exercícios};
			
			\node[
			text=eudark,
			font=\small
			]
			at (6.5,-0.45)
			{Física e Química A -- 10.º ano};
			
			\node[
			text=eudark,
			font=\small
			]
			at (6.5,-0.85)
			{Prof. José Gonçalves};
			
		\end{tikzpicture}
		
	\end{center}
	
\end{document}